\chapter{Practical Workflows}
\label{ch:workflows}

\index{workflows}

This chapter demonstrates how to use acme for common programming tasks.
The goal is not to prescribe a workflow but to show how acme's simple
tools compose into effective patterns.

\section{Editing Code}
\label{sec:coding}

\subsection{Opening a Project}

Start acme with the files you want to edit:

\begin{lstlisting}[language=bash]
acme src/*.c include/*.h Makefile
\end{lstlisting}

Or start acme empty and right-click (\button{3}) on directory names to
browse, then right-click on filenames to open them.

\subsection{Navigating}

\begin{itemize}
\item Right-click on a function name to search for it in the current
      file.
\item Right-click on \file{header.h} to open it.
\item Right-click on \texttt{:42} to jump to line 42.
\item Right-click on \texttt{:/pattern/} to search by regex.
\item Use \cmd{Look} with a 2-1 chord to search for specific text.
\end{itemize}

\subsection{Multiple Columns}

Use two or three columns to keep related files visible:

\begin{itemize}
\item Left column: header files and interfaces.
\item Middle column: implementation files you are editing.
\item Right column: \file{+Errors}, \cmd{Win} shell, or test output.
\end{itemize}

Create columns with \cmd{Newcol}.  Move windows between columns by
dragging the layout button.

\section{Building and Testing}
\label{sec:building}

\subsection{Running Make}

Type \cmd{make} in any tag bar and middle-click it.  Output appears in
\file{+Errors}.

If the build fails with errors like:

\begin{Verbatim}[frame=single]
main.c:42: error: undeclared identifier 'x'
\end{Verbatim}

\noindent
right-click on that line to jump directly to \file{main.c} line~42.
Fix the error, then middle-click \cmd{make} again.

\subsection{Running Tests}

\begin{lstlisting}[language=bash]
make test
go test ./...
python -m pytest
cargo test
\end{lstlisting}

Any of these can be typed in a tag and executed.  Test output with
file:line references is clickable.

\subsection{Continuous Workflow}

A typical edit-build-fix cycle:

\begin{enumerate}
\item Edit code in a window.
\item Middle-click \cmd{Put} to save.
\item Middle-click \cmd{make} in the tag.
\item Right-click error lines to navigate to problems.
\item Fix and repeat.
\end{enumerate}

For even faster iteration, open a \cmd{Win} window and run build
commands interactively.

\section{Version Control}
\label{sec:vcs}

\index{version control}
\index{git}

\subsection{Checking Status}

Type in any tag and middle-click:

\begin{lstlisting}[language=bash]
git status
git diff
git log --oneline -20
\end{lstlisting}

\subsection{Committing}

\begin{lstlisting}[language=bash]
git add -p
git commit
\end{lstlisting}

For interactive commands like \cmd{git add -p} or \cmd{git commit}
(which opens an editor), use a \cmd{Win} window.

\subsection{Viewing Diffs}

Run \cmd{git diff} and the output appears in \file{+Errors}.
Right-click on filenames in the diff headers to open those files.

\section{Writing Prose}
\label{sec:prose}

\subsection{Proportional Fonts}

Switch to a proportional font for more comfortable reading:

\begin{enumerate}
\item Middle-click \cmd{Font} in the tag to toggle between fixed and
      proportional.
\item Or specify a font: \cmd{Font /mnt/font/DejaVuSerif/14a/font}
\end{enumerate}

\subsection{Reflowing Text}

Select a paragraph and execute:

\begin{lstlisting}[language=bash]
|fmt -w 72
\end{lstlisting}

This reflows the text to 72 columns, suitable for email or plain-text
documents.

\subsection{Spell Checking}

Select text and run:

\begin{lstlisting}[language=bash]
>aspell list
\end{lstlisting}

This sends the selection to \cmd{aspell} and shows misspelled words in
\file{+Errors}.

\section{Multi-File Search and Replace}
\label{sec:multifile}

\subsection{Using Edit X}

To replace ``oldFunc'' with ``newFunc'' in all open \file{.c} files:

\begin{Verbatim}[frame=single]
Edit X/\.c$/ ,s/oldFunc/newFunc/g
\end{Verbatim}

\subsection{Using Grep}

Run \cmd{grep -rn pattern *.c} from a tag.  The output lines are
clickable:

\begin{Verbatim}[frame=single]
src/main.c:42:    oldFunc(x, y);
src/util.c:17:    oldFunc(a, b);
\end{Verbatim}

Right-click each line to jump to the match, make the change, and move
to the next.

\section{Session Management}
\label{sec:sessions}

\index{session management}
\index{Dump}
\index{Load}

\subsection{Saving Your Workspace}

Before quitting, middle-click \cmd{Dump} in the top tag.  This saves
the complete state---all windows, their positions, fonts, and column
layout---to \file{\$HOME/acme.dump}.

\subsection{Restoring}

Next time you start acme:

\begin{lstlisting}[language=bash]
acme -l $HOME/acme.dump
\end{lstlisting}

Or start acme normally and middle-click \cmd{Load} in the top tag.

\subsection{Multiple Sessions}

Save different sessions to different files:

\begin{Verbatim}[frame=single]
Dump project-a.dump
Dump project-b.dump
\end{Verbatim}

\noindent
Then load the appropriate one:

\begin{Verbatim}[frame=single]
Load project-a.dump
\end{Verbatim}

\section{Tips}

\begin{itemize}
\item Add frequently-used commands to the top tag bar.  For example,
      type \cmd{make} or \cmd{git diff} there so it is always one
      click away.
\item Use directory windows as a file browser---right-click a directory
      path to list it.
\item Right-clicking on a word in \file{+Errors} output searches for
      it, which is useful for tracing error messages.
\item The \cmd{Zerox} command creates a second view of a file, handy for
      looking at one part while editing another.
\item \cmd{Sort} in a column tag alphabetizes windows, making it easy
      to find files.
\end{itemize}
